\section{容斥原理与重复计数}
\label{sec:04-Discrete-Mathematics/01-Combinatorics/04}

一家初创公司有 300 个注册用户。运营团队拿出了三份触达数据：150 人收到了应用推送，120 人打开了邮件，100 人参加了线下活动。老板问：总共有多少不同的人被触达？你提议直接加：\(150+120+100=370\)。老板看着报表上的``300 注册用户''皱了皱眉。

出问题的地方一目了然：有人同时出现在了多个渠道里，被重复数了。把三个渠道的人数简单相加，每个跨渠道的用户都被多算了两到三次。你需要一个方法，把多算的扣掉，把扣多了的补回来。

\subsection{直接相加为何会多算交集}

先看两个集合的情况。设 \(A\) 是收到推送的人，\(B\) 是打开邮件的人。\(|A|+|B|\) 把同时属于两边的用户数了两次。把多算的那一次减去：

\[
|A\cup B|
=|A|+|B|-|A\cap B|.
\]

三个集合时，事情开始变得绕。你先加单集合：\(|A|+|B|+|C|\)。这时候每个恰好属于两个集合的人被数了两次。于是你减去所有两两交集：\(|A\cap B|+|A\cap C|+|B\cap C|\)。但等等——如果有个人三个渠道都用了呢？在加单集合时他被数了 3 次，减两两交集时他被减了 3 次——现在他的计数是 \(3-3=0\)，消失了。这就是被``减过头''了。你得把三重交集加回来：

\[
|A\cup B\cup C|
=|A|+|B|+|C|-|A\cap B|-|A\cap C|-|B\cap C|+|A\cap B\cap C|.
\]

这个结构已经能看出交替符号的苗头：加单数、减两两、加三三。但为什么恰好是这个模式，而不是别的什么修正方案？

\subsection{每个并集元素最终权重恰好为 1}

让我们追踪一个人的命运。假设这个人恰好出现在了 \(r\) 个渠道中（\(r\ge 1\)）。

在``加单集合''那一层，他所属的 \(r\) 个集合各被加一次，他被数了 \(\binom r1\) 次。

在``减两两交集''那一层，从他的 \(r\) 个渠道中任选两个，那个交集都包含他。他被减了 \(\binom r2\) 次。

到了``加三三交集''，他从 \(r\) 个渠道中任选三个，对应交集同样包含他。他被加回 \(\binom r3\) 次。

继续这个模式，他在整个容斥过程中最终的净权重是

\[
\binom r1-\binom r2+\binom r3-\cdots+(-1)^{r+1}\binom rr.
\]

这个和等于什么？回忆二项式展开 \((1-1)^r=\sum_{j=0}^r\binom rj 1^{r-j}(-1)^j\)。展开后 \(\binom r0=1\)，后面的交替和恰好是我们上面写的那个式子（只是符号方向差一个整体负号）。精确地说：

\[
0=(1-1)^r=\binom r0-\binom r1+\binom r2-\cdots+(-1)^r\binom rr,
\]

所以

\[
\binom r1-\binom r2+\cdots+(-1)^{r+1}\binom rr
=1-\bigl(1-1\bigr)^r=1.
\]

每一个属于至少一个集合的元素，最终被精确地数了一次——不重、不漏。容斥不是什么神秘的``交替加减口诀''，它就是一套逐层修正的权重方案：先承认交集对象被多算，然后让交替的二项式系数恰好把多算的部分抵消到只剩 1。

回到那个 300 个用户的例子。容斥能纠正重复计数，但有一个前提：同一个人必须能在不同渠道间被识别为同一个集合元素。如果一位用户在推送系统中用的手机号和邮件系统中用的邮箱没有关联起来，交集就会被低估，并集仍然被高估。容斥修的是计数逻辑的重叠，不能替数据修复错误的身份对应。这个区分在工程中容易被忽略：公式是正确的，但喂给公式的数据可能已经漏了或者重了。

\subsection{一个整除计数例子}

1 到 100 中，能被 2 或 3 整除的数有多少个？

把范围里的整数当成集合。设 \(A\) 是 2 的倍数，\(B\) 是 3 的倍数。

能被 2 整除：\(\lfloor 100/2\rfloor=50\) 个。
能被 3 整除：\(\lfloor 100/3\rfloor=33\) 个。
同时被 2 和 3 整除，就是被最小公倍数 6 整除：\(\lfloor 100/6\rfloor=16\) 个。

容斥给出：

\[
|A\cup B|=50+33-16=67.
\]

即 1 到 100 中有 67 个数能被 2 或 3 整除。如果想问``既不被 2 也不被 3 整除''，用全集 100 减去并集：\(100-67=33\)。

这个例子用两个集合就做完了，但思路可以自然地扩展到更多除数。比如说，能被 2、3 或 5 整除的数有多少个？三集合容斥，依次加单数、减两两交集、加三三交集。补集法常常把``至少一个''或``没有任何''这类条件翻译得更简单；容斥则负责纠正多个``坏事件''之间的重叠。

\subsection{错排问题}

\(n\) 封信，\(n\) 个对应的信封。把信随机放进去，要求没有一封放对——这就是错排问题。

令 \(A_i\) 表示第 \(i\) 封信放对了这个事件。所有排列共 \(n!\) 个，其中至少一封放对的排列数是 \(|A_1\cup\cdots\cup A_n|\)。错排数 \(D_n\) 是全集减去这个并集：

\[
D_n=n!-|A_1\cup\cdots\cup A_n|.
\]

固定某 \(k\) 封信让它们放对——选哪 \(k\) 封有 \(\binom nk\) 种，选好后剩下的 \(n-k\) 封可以任意排列，有 \((n-k)!\) 种。所以一个 \(k\) 重交集的大小是 \(\binom nk(n-k)!\)。容斥公式里，这项带着符号 \((-1)^k\)。注意 \(\binom nk(n-k)!=n!/k!\)，于是

\[
D_n
=n!-\sum_{k=1}^n(-1)^{k-1}\frac{n!}{k!}
=n!\sum_{k=0}^n\frac{(-1)^k}{k!}.
\]

算出几个小 \(n\) 来感受一下：\(D_1=0\)（就一封信，不放对是不可能的），\(D_2=1\)（两封信互换），\(D_3=2\)，\(D_4=9\)，\(D_5=44\)。

这个公式还藏着另一个信息：当 \(n\) 很大时，括号里的和逼近 \(e^{-1}\) 的泰勒展开。所以 \(D_n\approx n!/e\)——错排数大约是全部排列数的 \(1/e\approx 0.368\)。容斥把一个``每个位置都禁止''的全局条件，拆成了若干``指定哪些位置必须正确''的易计数的局部条件，然后交替修正重叠。这就是容斥最擅长的翻译工作。

\subsection{截断容斥给上下界}

集合很多的时候，完整容斥公式的项数是指数级的——\(n\) 个集合有 \(2^n-1\) 个非空交集。全部算完可能不现实。

这时候，只取前几层就能得到实用的上下界。这个截断版本叫 Bonferroni 不等式：

\begin{itemize}
\tightlist
\item
  只加单集合的和：给出并集的上界（你多算了，所以偏大）；
\item
  加到减去两两交集：给出并集的下界（你可能减过头了，所以偏小）；
\item
  再加回三三交集：又回到上界；
\item
  继续交替。
\end{itemize}

概率论里的 union bound 正是只保留第一层的情况：

\[
P\Bigl(\bigcup_i A_i\Bigr)\le\sum_i P(A_i).
\]

这个上界虽然粗糙——它完全忽略事件之间的重叠——但不需要知道任何交集的概率，因此适用范围极广。当你面对几十个可能相互纠缠的坏事件，而又算不出它们的联合分布时，union bound 常常是唯一可行的保守估计。粗糙不要紧，能用就行。
